% profiles/university-of-minnesota-twin-cities.tex — University of Minnesota-Twin Cities
% ============================================================================
% source: https://onestop.umn.edu/academics/grad-and-professional/thesis-or-dissertation-submission
%         (the live submission page, which says "You must follow the University
%          of Minnesota formatting and submission guidelines" and links the PDF
%          below; it also warns that "The guidelines found in your department or
%          in the University Library to format your thesis may not meet current
%          formatting guidelines")
% source: https://drive.google.com/file/d/1-EFscAFvtZbJXyPkCs4OgPQ6IHrAJjEj/view
%         ("Thesis Formatting & Submission Guidelines", Graduate Student
%          Services and Progress; every page footed "Updated: January 2023")
% accessed: 2026-09-03
% checked:  texlib-thesis-profile-round
%
% VERIFIED and set here: geometry, spacing, the title page, the absence of an
% approval page, and the absence of any published accessibility requirement.
% NOT set (left to the class's neutral behaviour):
%   * \thesissetchapteropening — Minnesota states no deeper margin on chapter
%     openings. The margin rule is flat: "Every page of the thesis, including
%     all appendices, notes, and the bibliography".
%   * a font rule — Minnesota's font list is headed "Recommended font types",
%     and "The recommended minimum font size is 10-point", so there is nothing
%     to require. The class's Latin Modern satisfies "Standard fonts should be
%     used" and embeds every face.
%
% IPEDS records this institution as "University of Minnesota-Twin Cities"; the
% Graduate School, the guidelines and the prescribed title page all say
% "University of Minnesota", and that is what \thesisinstitution carries, since
% it is the string the class prints in warnings.
%
% CURRENCY. The governing document is dated January 2023 — three and a half
% years old at the access date. A second, OLDER edition is still served at
% https://assets.asr.umn.edu/files/gssp/Thesis_formatting_guidelines.pdf
% (10 pages; its PDF metadata says Adobe InDesign CS5, created 2013,
% modified 2021, and it carries no "Updated:" footer). It is NOT linked from
% the One Stop page, and it disagrees on the point that matters most here: it
% prescribes an in-document signature page reading "This is to certify that I
% have examined this copy of a [Master's thesis] or [Doctoral dissertation]
% by …", with a rule for the adviser's signature. The January 2023 edition
% replaced that with the online Approval and Deposit Agreement. This profile
% follows the January 2023 edition because One Stop links it and says it
% governs. If a reviewer finds that GSSP has since reinstated the signature
% leaf, the line to change is \thesisnoapprovalpage at the bottom.

\thesisinstitution{University of Minnesota}

% "At minimum, students must use 1 inch (2.6 cm) margins on all sides", applying
% to "Every page of the thesis, including all appendices, notes, and the
% bibliography", and "Nothing may appear in the margins, meaning no page
% numbers, text, tables, graphs, charts, or parts of illustrations."
%
% The 1.5in left margin is a recommendation and only for a printed copy — "It is
% also recommended that students consider using a 1.5 inch (3.9 cm) left margin
% if they intend to have their manuscript printed" — so it does not belong in
% the filed PDF's geometry. Add a binding offset at print time.
\thesissetgeometry{left=1in, right=1in, top=1in, bottom=1in}

% "The body of the thesis must be double-spaced or 1½-spaced." Minnesota offers
% the choice; this profile takes double, which is also the class default.
% \thesissetspacing{onehalf} is equally in spec.
%
% The rule is about the BODY. Single spacing is permitted for several parts and
% the class does not force those either way: "Long quotations, notes, and the
% bibliography may be single-spaced, unless the student's graduate program
% requires otherwise", "Captions may be single-spaced", and, for appendices,
% "text and information in appendices can be single-spaced".
\thesissetspacing{double}

% --- Accessibility -----------------------------------------------------------
% Read and found ABSENT, not skipped. The One Stop submission page and all eight
% pages of the January 2023 guidelines were read in full: neither states an
% accessibility requirement for the filed document. The word "accessible"
% appears in them only about the Digital Conservancy making the deposit
% "freely accessible online to the public". The requirements Minnesota does put
% on the file are format and size, not conformance.
%
% Nothing is therefore declared with \thesisrequiretagging,
% \thesisrequiredocumentlanguage or \thesisrequirepdfstandard. This class emits
% a tagged PDF regardless; declaring a requirement GSSP has not published would
% put words in their mouth.
%
% Elsewhere in the University — the Libraries' Digital Conservancy deposit
% guidance, the Accessible U site, and University policy on digital content —
% accessibility guidance does exist. None of it is the Graduate School's filing
% requirement, so none of it is recorded here.
\thesisaccessibilityrequirement{Not published}
	{Graduate Student Services and Progress states no accessibility requirement
	 for a filed thesis or dissertation. Its requirements for the file are that
	 "The full text must be in Adobe PDF format and must be one file" and that
	 the complete file size is no more than 1,000 MB. Read 2026-09-03 in the
	 January 2023 Formatting \& Submission Guidelines and on the One Stop
	 submission page; an absence found, not a section skipped.}

% --- Title page --------------------------------------------------------------
% The guidelines print a specimen title page ("See Title Page example below:")
% and the wording and line breaks here are that specimen, reproduced. The
% surrounding rules are:
%   1. "The title page is included as the first page."
%   2. "The title page is not to be numbered or counted."
%   3. "The year that degree requirements were met must be included (not
%      necessarily the year the student defends)."
%   5. "The title page must include only one title."
%   6. "The student's name listed on the title page may be the official name,
%      preferred name, or degree name on record with the University."
% and, from the checklist, "The title page has the correct year listed. Year
% must correspond to the year the degree is awarded."
%
% \graddate should therefore be given the YEAR of conferral. The class default
% is a month and year; a month is not forbidden, but the year is the part
% Minnesota checks, so whatever is passed must contain the right one.
%
% Rule 4 — "The title of the thesis must not contain chemical or mathematical
% formulas, symbols, superscripts, subscripts, Greek letters, or other
% non-standard characters; words must be substituted" — is a constraint on
% \title, which no macro can enforce. It is stated here so a mathematician
% reads it before typing $\Gamma$ into the title.
%
% The specimen sets the boilerplate in capitals and leaves the title in title
% case; the degree line is capitalised there too, so \thesis@degree is
% uppercased. VERTICAL placement is not prescribed anywhere in the document —
% only the wording, the order and the line breaks are — so the gaps below are a
% reading of the specimen's proportions, not a quoted figure.
\renewcommand{\thesistitlepage}{%
	\loadgeometry{nopagenumber}%
	\begin{titlepage}%
		\thispagestyle{empty}\singlespacing\centering
		\vspace*{1in}%
		\thesis@tagtitle{Title}{\@title}\par
		\vspace{0.75in}
		A \thesis@DocNoun\\
		SUBMITTED TO THE FACULTY OF THE UNIVERSITY\\
		OF MINNESOTA\\
		BY\par
		\vspace{0.75in}
		\@author\par
		\vspace{0.75in}
		IN PARTIAL FULFILLMENT OF THE REQUIREMENTS\\
		FOR THE DEGREE OF\\
		\MakeUppercase{\thesis@degree}\par
		\vspace{0.75in}
		\thesis@advisor\par
		\vspace{0.75in}
		\thesis@date\par
		\vspace*{\fill}
	\end{titlepage}%
	\loadgeometry{normal}%
}

% --- Committee approval page --------------------------------------------------
% Minnesota prescribes none inside the document. The introductory material is
% enumerated in the guidelines — title page, copyright page, acknowledgements,
% dedication, abstract, table of contents, list of tables, list of figures, then
% "other introductory material (e.g., notations, list of abbreviations)" — and
% no approval or signature leaf is among them. Approval is collected on a
% separate online form instead:
%
%   "Students must sign and submit the online form for Thesis/Dissertation
%    Approval and Deposit Agreement … It also includes the signature page, which
%    is signed by the advisor (and co-advisor if applicable). The advisor(s)
%    will confirm they have seen and approve of the final version of the
%    thesis."
%
% So the signature page exists, but it is part of that form and not part of the
% filed PDF, which "must be one file" containing "all introductory pages, the
% body of the manuscript, bibliography, and appendices". The adviser is named on
% the title page above.
%
% Note what is NOT quoted: Minnesota nowhere writes that the document may not
% contain an approval page, the way Washington does. The page is taken away on
% the strength of the enumerated introductory material and the separate signed
% form, and because the superseded 2013/2021 edition's signature leaf is exactly
% the page a student would otherwise carry over. See the CURRENCY note above.
\thesisnoapprovalpage
